\documentclass[12pt]{article}
\usepackage{amsmath, amssymb, amsthm}

\title{Superposition Principle, Thévenin's Theorem, and Norton's Theorem}
\author{}
\date{}

\begin{document}
	\maketitle
	
	\section*{0. Superposition Principle (Purely Mathematical)}
	
	Consider a linear system
	$$
	\mathbf{A}\mathbf{x} = \mathbf{b},
	$$
	where $\mathbf{A} \in \mathbb{R}^{m \times m}$ is nonsingular, 
	$\mathbf{x} \in \mathbb{R}^m$ is the unknown vector, 
	and $\mathbf{b} \in \mathbb{R}^m$ is the excitation vector.  
	
	Suppose $\mathbf{b}$ can be decomposed as
	$$
	\mathbf{b} = \mathbf{b}_1 + \mathbf{b}_2 + \cdots + \mathbf{b}_n.
	$$
	
	By linearity of matrix equations,
	$$
	\mathbf{x} = \mathbf{A}^{-1}\mathbf{b} 
	= \mathbf{A}^{-1}\mathbf{b}_1 + \mathbf{A}^{-1}\mathbf{b}_2 + \cdots + \mathbf{A}^{-1}\mathbf{b}_n.
	$$
	
	Defining $\mathbf{x}_k = \mathbf{A}^{-1}\mathbf{b}_k$, we obtain
	$$
	\mathbf{x} = \sum_{k=1}^n \mathbf{x}_k.
	$$
	
	This is the \textbf{superposition principle}: the total response of a linear system is the sum of the responses to each input considered separately.  
	
	---
	
	\section*{1. Thévenin's Theorem}
	
	\subsection*{Formal Statement}
	
	Let a linear network $\mathcal{N}$ with terminals $A,B$ consist of independent sources and linear elements.  
	For any external load $R_L$ connected at $A,B$, the network is equivalent to
	$$
	\mathcal{N}_{AB} \equiv V_{\text{th}} \;\; \text{in series with} \;\; R_{\text{th}},
	$$
	where $V_{\text{th}}$ is the open-circuit voltage and $R_{\text{th}}$ the Thévenin resistance.  
	
	\subsection*{Derivation Basis}
	
	Because the circuit is linear, the terminal voltage must be an affine function of the load current $I_L$:
	$$
	V_{AB} = V_{\text{th}} - R_{\text{th}} I_L.
	$$
	
	This affine relation follows directly from superposition and the linearity of Kirchhoff's equations.
	
	\subsection*{Computational Formulae}
	
	1. Open-circuit voltage:
	$$
	V_{\text{th}} = V_{AB}\Big|_{I_L = 0}.
	$$
	
	2. Thévenin resistance (two methods):
	
	(a) Deactivation method:
	$$
	R_{\text{th}} = R_{\text{eq}}(A,B) \quad \text{with independent sources set to zero}.
	$$
	
	(b) Short-circuit method:
	$$
	R_{\text{th}} = \frac{V_{\text{th}}}{I_{\text{sc}}}.
	$$
	
	3. Load response:
	$$
	I_L = \frac{V_{\text{th}}}{R_{\text{th}} + R_L}, 
	\qquad
	V_L = \frac{R_L}{R_{\text{th}} + R_L} V_{\text{th}}.
	$$
	
	\subsection*{Special Case: Maximum Power Transfer}
	
	The power on the load is
	$$
	P_L = \frac{V_{\text{th}}^2}{(R_{\text{th}} + R_L)^2} \, R_L.
	$$
	
	This is maximized when
	$$
	R_L = R_{\text{th}},
	$$
	with maximum power
	$$
	P_{L,\text{max}} = \frac{V_{\text{th}}^2}{4R_{\text{th}}}.
	$$
	
	---
	
	\section*{2. Norton's Theorem}
	
	\subsection*{Formal Statement}
	
	Let a linear network $\mathcal{N}$ with terminals $A,B$ consist of independent sources and linear elements.  
	Then, for any load $R_L$ connected across $A,B$, the network is equivalent to
	$$
	\mathcal{N}_{AB} \equiv I_{\text{N}} \;\; \text{in parallel with} \;\; R_{\text{N}},
	$$
	where $I_{\text{N}}$ is the short-circuit current at the terminals and $R_{\text{N}}$ the Norton resistance.
	
	\subsection*{Algebraic Relation to Thévenin}
	
	From Thévenin’s theorem we have
	$$
	V_{AB} = V_{\text{th}} - R_{\text{th}} I_L.
	$$
	
	If the terminals are short-circuited, $V_{AB} = 0$, so
	$$
	I_{\text{sc}} = \frac{V_{\text{th}}}{R_{\text{th}}}.
	$$
	
	Define
	$$
	I_{\text{N}} = I_{\text{sc}}, 
	\qquad 
	R_{\text{N}} = R_{\text{th}}.
	$$
	
	Thus, the Norton form is
	$$
	\mathcal{N}_{AB} \equiv I_{\text{N}} \;\; \text{in parallel with} \;\; R_{\text{N}}.
	$$
	
	\subsection*{Load Response in Norton Form}
	
	For a load $R_L$ connected in parallel:
	$$
	I_L = \frac{I_{\text{N}} R_{\text{N}}}{R_{\text{N}} + R_L}, 
	\qquad
	V_L = \frac{R_L}{R_{\text{N}} + R_L} I_{\text{N}} R_{\text{N}}.
	$$
	
	\subsection*{Equivalence of Forms}
	
	Since
	$$
	V_{\text{th}} = I_{\text{N}} R_{\text{N}}, \qquad R_{\text{th}} = R_{\text{N}},
	$$
	the Thévenin and Norton models are algebraically dual and interchangeable.
	
	\section*{Example: Superposition Principle (Pure Algebraic Case)}
	
	Consider the linear system
	$$
	\begin{cases}
		2x + y = b_1 + b_2, \\
		x - y = c_1 + c_2.
	\end{cases}
	$$
	
	Here the right-hand side is the sum of two independent inputs 
	$(b_1, c_1)$ and $(b_2, c_2)$.
	
	---
	
	\subsection*{Step 1. Decompose System}
	
	We split into two subsystems:
	
	1. With input $(b_1, c_1)$:
	$$
	\begin{cases}
		2x + y = b_1, \\
		x - y = c_1,
	\end{cases}
	$$
	solution $(x_1, y_1)$.
	
	2. With input $(b_2, c_2)$:
	$$
	\begin{cases}
		2x + y = b_2, \\
		x - y = c_2,
	\end{cases}
	$$
	solution $(x_2, y_2)$.
	
	---
	
	\subsection*{Step 2. Solve Each Subsystem}
	
	From $x - y = c_1$ we have $x = y + c_1$. Substituting into $2x + y = b_1$:
	$$
	2(y+c_1) + y = b_1 \;\;\implies\;\; 3y + 2c_1 = b_1.
	$$
	Thus
	$$
	y_1 = \frac{b_1 - 2c_1}{3}, 
	\qquad
	x_1 = \frac{b_1 - 2c_1}{3} + c_1.
	$$
	
	Similarly, for $(b_2, c_2)$:
	$$
	y_2 = \frac{b_2 - 2c_2}{3}, 
	\qquad
	x_2 = \frac{b_2 - 2c_2}{3} + c_2.
	$$
	
	---
	
	\subsection*{Step 3. Apply Superposition}
	
	By the superposition principle,
	$$
	x = x_1 + x_2, 
	\qquad 
	y = y_1 + y_2.
	$$
	
	That is,
	$$
	x = \left(\frac{b_1 - 2c_1}{3} + c_1\right) 
	+ \left(\frac{b_2 - 2c_2}{3} + c_2\right),
	$$
	$$
	y = \frac{b_1 - 2c_1}{3} + \frac{b_2 - 2c_2}{3}.
	$$
	
	---
	
	\subsection*{Step 4. Verification}
	
	Solving directly with combined input $(b_1+b_2,\,c_1+c_2)$ gives
	$$
	y = \frac{(b_1+b_2) - 2(c_1+c_2)}{3}, 
	\qquad
	x = y + (c_1+c_2),
	$$
	which agrees with the sum $(x_1+x_2,\, y_1+y_2)$. \qed
	 
	
\end{document}
